\documentclass[aspectratio=169,xcolor=dvipsnames]{beamer}
\usetheme{Berlin}

\usepackage[spanish]{babel}
\usepackage{hyperref}
\usepackage{graphicx}
\usepackage{booktabs}
\usepackage{amsmath}
\usepackage{lettrine}
\setbeamertemplate{caption}[numbered]
\usepackage[dvipsnames,svgnames,x11names]{xcolor}
\usepackage{xurl}
\usepackage{algorithm}
\usepackage{algorithmicx}
\usepackage{algpseudocode}
\usepackage{adjustbox}

\hypersetup{
    colorlinks=true,
    linkcolor=cyan,
    filecolor=blue,
    urlcolor=blue,
    citecolor=blue,
}

%----------------------------------------------------------------------------------------
%   CODE LISTINGS SETTINGS
%----------------------------------------------------------------------------------------
\usepackage{listings}
\definecolor{backcolour}{rgb}{0.97,0.97,0.99}
\definecolor{codegreen}{rgb}{0,0.6,0}

\lstdefinestyle{MATLABStyle}{
  language=Matlab,
  basicstyle=\ttfamily\scriptsize,
  keywordstyle=\color{blue}\bfseries,
  commentstyle=\color{codegreen},
  stringstyle=\color{violet},
  breaklines=true,
  numbers=left,
  numbersep=5pt,
  frame=lines,
  backgroundcolor=\color{backcolour}
}
\lstset{style=MATLABStyle}

%----------------------------------------------------------------------------------------
%   TITLE CONFIGURATION
%----------------------------------------------------------------------------------------
\title{Simulated Annealing}
\subtitle{Metaheurística basada en el enfriamiento de materiales}
\author{Prof. Dr. BARSEKH-ONJI Aboud}
\institute{
    Facultad de Ingeniería \\
    Universidad Anáhuac México
}
\date{\today}

%----------------------------------------------------------------------------------------
%   AUTO AGENDA AT EACH SECTION
%----------------------------------------------------------------------------------------
\AtBeginSection[]{
  \begin{frame}{Agenda}
    \tableofcontents[currentsection]
  \end{frame}
}

\begin{document}

\begin{frame}
    \titlepage
\end{frame}

\begin{frame}{Contenido General}
    \tableofcontents
\end{frame}

%------------------------------------------------
\section{Motivación y Contexto}
%------------------------------------------------

\begin{frame}{¿Por qué necesitamos metaheurísticas?}
    \begin{block}{El problema de la optimización compleja}
        Muchos problemas de ingeniería presentan:
        \begin{itemize}
            \item Espacios de búsqueda \textbf{no convexos} con múltiples mínimos locales
            \item Funciones objetivo \textbf{no diferenciables} o de cómputo costoso
            \item Alta dimensionalidad que hace inviables los métodos exactos
        \end{itemize}
    \end{block}
    \vspace{0.3cm}
    \begin{alertblock}{Limitación de métodos clásicos}
        El descenso por gradiente y otros métodos locales \textbf{quedan atrapados} en mínimos locales. No pueden explorar globalmente el espacio de soluciones.
    \end{alertblock}
\end{frame}

\begin{frame}{Inspiración de la Naturaleza}
    \begin{columns}
        \column{0.55\textwidth}
        \begin{block}{Analogía con la metalurgia}
            El \textbf{\textit{recocido}} (\textit{annealing}) es un proceso térmico donde:
            \begin{enumerate}
                \item El material se calienta a alta temperatura
                \item Los átomos adquieren \textbf{alta energía} y se reorganizan libremente
                \item Se enfría \textbf{lentamente y de forma controlada}
                \item Los átomos llegan a un estado de mínima energía (estructura cristalina óptima)
            \end{enumerate}
        \end{block}
        \column{0.45\textwidth}
        \begin{exampleblock}{Origen del algoritmo}
            Propuesto por \textbf{Kirkpatrick, Gelatt y Vecchi} (1983) en el artículo:
            \textit{``Optimization by Simulated Annealing''}, \textit{Science}, Vol. 220.
            \vspace{0.2cm}
            
            Inspirado en la simulación de Monte Carlo de Metrópolis (1953).
        \end{exampleblock}
    \end{columns}
\end{frame}

%------------------------------------------------
\section{Fundamentos Teóricos}
%------------------------------------------------

\begin{frame}{El Criterio de Metrópolis}
    El núcleo matemático del algoritmo se basa en la \textbf{distribución de Boltzmann}:
    \begin{equation}
        P(\Delta E, T) = \exp\!\left(\frac{-\Delta E}{T}\right)
    \end{equation}
    \begin{itemize}
        \item $\Delta E = f(x_{\text{nuevo}}) - f(x_{\text{actual}})$: diferencia de costo entre soluciones
        \item $T$: temperatura actual del sistema (parámetro de control)
        \item $P$: probabilidad de \textbf{aceptar una solución peor}
    \end{itemize}
    \vspace{0.2cm}
    \begin{alertblock}{Propiedad clave}
        Cuando $T$ es \textbf{alto}, $P \to 1$: se aceptan soluciones malas con frecuencia (exploración). \\
        Cuando $T$ es \textbf{bajo}, $P \to 0$: solo se aceptan mejoras (explotación).
    \end{alertblock}
\end{frame}

\begin{frame}{Esquema de Enfriamiento}
    La temperatura disminuye iterativamente mediante un \textbf{factor de enfriamiento} $\alpha$:
    \begin{equation}
        T_{k+1} = \alpha \cdot T_k, \qquad 0 < \alpha < 1
    \end{equation}
    \begin{columns}
        \column{0.5\textwidth}
        \begin{block}{Parámetros de control}
            \begin{itemize}
                \item $T_0$: Temperatura inicial (alta)
                \item $T_{\min}$: Temperatura mínima (criterio de paro)
                \item $\alpha$: Factor de enfriamiento ($0.90$--$0.99$)
                \item $N_{\max}$: Iteraciones máximas
            \end{itemize}
        \end{block}
        \column{0.5\textwidth}
        \begin{alertblock}{Impacto de $\alpha$}
            \begin{itemize}
                \item $\alpha$ pequeño ($< 0.9$): enfriamiento rápido, riesgo de mínimo local
                \item $\alpha$ grande ($> 0.99$): enfriamiento lento, mejor exploración pero mayor costo computacional
            \end{itemize}
        \end{alertblock}
    \end{columns}
\end{frame}

\begin{frame}{Balance entre Exploración y Explotación}
    \begin{columns}
        \column{0.5\textwidth}
        \textbf{Exploración} (\textit{Exploration}):
        \begin{itemize}
            \item Temperaturas altas ($T$ grande)
            \item Se aceptan soluciones peores
            \item El algoritmo ``escapa'' de mínimos locales
            \item Mayor diversidad en el espacio de búsqueda
        \end{itemize}
        \column{0.5\textwidth}
        \textbf{Explotación} (\textit{Exploitation}):
        \begin{itemize}
            \item Temperaturas bajas ($T$ pequeño)
            \item Solo se aceptan mejoras
            \item Convergencia hacia el óptimo
            \item Refinamiento de la solución actual
        \end{itemize}
    \end{columns}
    \vspace{0.4cm}
    \begin{exampleblock}{Principio fundamental}
        El éxito de SA reside en \textbf{combinar} ambas fases de forma adaptativa mediante la temperatura. El proceso imita la ``paciencia térmica'' del recocido metalúrgico.
    \end{exampleblock}
\end{frame}

%------------------------------------------------
\section{Pseudocódigo y Algoritmo}
%------------------------------------------------

\begin{frame}{Pseudocódigo del Simulated Annealing}
\begin{adjustbox}{max width=0.97\textwidth}
\begin{minipage}{\textwidth}
\begin{algorithm}[H]
\tiny
\caption{Simulated Annealing}
\begin{algorithmic}[1]
    \State \textbf{Inicializar:} $x_0$ (solución inicial), $T_0$, $T_{\min}$, $\alpha$, $N_{\max}$
    \State $x_{\text{mejor}} \leftarrow x_0$; \quad $f_{\text{mejor}} \leftarrow f(x_0)$
    \State $T \leftarrow T_0$; \quad $k \leftarrow 0$
    \While{$T > T_{\min}$ \textbf{and} $k < N_{\max}$}
        \State Generar vecino: $x_{\text{nuevo}} \leftarrow x_k + \mathcal{N}(0,1)$
        \State Calcular $\Delta E = f(x_{\text{nuevo}}) - f(x_k)$
        \If{$\Delta E < 0$}
            \State $x_{k+1} \leftarrow x_{\text{nuevo}}$ \Comment{Aceptar mejora}
            \If{$f(x_{\text{nuevo}}) < f_{\text{mejor}}$}
                \State $x_{\text{mejor}} \leftarrow x_{\text{nuevo}}$
            \EndIf
        \Else
            \If{$\text{rand}(0,1) < \exp(-\Delta E / T)$}
                \State $x_{k+1} \leftarrow x_{\text{nuevo}}$ \Comment{Aceptar con probabilidad}
            \EndIf
        \EndIf
        \State $T \leftarrow \alpha \cdot T$; \quad $k \leftarrow k + 1$
    \EndWhile
    \State \Return $x_{\text{mejor}}$
\end{algorithmic}
\end{algorithm}
\end{minipage}
\end{adjustbox}
\end{frame}

%------------------------------------------------
\section{Función de Prueba: Ackley}
%------------------------------------------------

\begin{frame}{La Función de Ackley}
    Una de las funciones de referencia más utilizadas en optimización:
    \begin{equation}
        f(x,y) = -20\exp\!\left(-0.2\sqrt{0.5(x^2+y^2)}\right)
        - \exp\!\left(0.5(\cos 2\pi x + \cos 2\pi y)\right) + 20 + e
    \end{equation}
    \begin{columns}
        \column{0.55\textwidth}
        \begin{block}{Características}
            \begin{itemize}
                \item \textbf{Mínimo global:} $f(0,0) = 0$
                \item \textbf{Dominio típico:} $[-5, 5]^2$
                \item \textbf{Superficie:} irregular, con múltiples mínimos locales
                \item Diseñada para \textbf{engañar} a algoritmos que solo explotan localmente
            \end{itemize}
        \end{block}
        \column{0.45\textwidth}
        \begin{alertblock}{Desafío para optimizadores}
            La función tiene una región central casi plana rodeada de ``muros'' cosénicos. Un optimizador local queda atrapado fuera del óptimo global.
        \end{alertblock}
    \end{columns}
\end{frame}

\begin{frame}[fragile]{Implementación MATLAB — Definición de la Función}
\begin{columns}
\column{0.6\textwidth}
\begin{exampleblock}{Función anónima en MATLAB}
Definición compacta usando \texttt{@(x)} para pasar un vector $[x_1, x_2]$:
\end{exampleblock}
\column{0.4\textwidth}
\begin{block}{Parámetros matemáticos}
\begin{itemize}
    \item $a = 20$, $b = 0.2$, $c = 2\pi$
    \item Estándar de la literatura
\end{itemize}
\end{block}
\end{columns}
\vspace{0.2cm}
\begin{lstlisting}[language=Matlab, style=MATLABStyle]
% Definicion de la funcion de Ackley como funcion anonima
ackley = @(x) -20 * exp(-0.2 * sqrt(0.5 * (x(1)^2 + x(2)^2))) ...
              - exp(0.5 * (cos(2 * pi * x(1)) + cos(2 * pi * x(2)))) ...
              + 20 + exp(1);

% Generacion de la malla para visualizacion
[x_grid, y_grid] = meshgrid(-5:0.1:5, -5:0.1:5);
z_grid = arrayfun(@(x, y) ackley([x, y]), x_grid, y_grid);
\end{lstlisting}
\end{frame}

%------------------------------------------------
\section{Implementación MATLAB}
%------------------------------------------------

\begin{frame}[fragile]{Parámetros e Inicialización}
\begin{lstlisting}[language=Matlab, style=MATLABStyle]
% ---- Parametros del algoritmo ----
T       = 1000;   % Temperatura inicial (alta -> maxima exploracion)
T_min   = 1;      % Temperatura minima (criterio de paro)
alpha   = 0.95;   % Factor de enfriamiento geometrico
max_iter = 700;   % Numero maximo de iteraciones

% ---- Solucion inicial ----
x = [0.1, 0.1];          % Punto de partida (cercano al origen)
best_sol  = x;
best_cost = ackley(x);   % Evaluar costo inicial

% ---- Vectores de rastreo ----
costs   = zeros(1, max_iter);
trace_x = x(1);  trace_y = x(2);  trace_z = best_cost;
\end{lstlisting}
\end{frame}
\begin{frame}[fragile]{Parámetros e Inicialización}

\begin{block}{Nota sobre $\alpha = 0.95$}
    Con $T_0=1000$, $T_{\min}=1$ y $\alpha=0.95$: el número de iteraciones para enfriar es $\lceil \log(T_{\min}/T_0)/\log(\alpha) \rceil \approx 135$ pasos de temperatura.
\end{block}
\end{frame}

\begin{frame}[fragile]{Lazo Principal del SA}
\begin{lstlisting}[language=Matlab, style=MATLABStyle]
for iter = 1:max_iter
    % 1. Generar solucion vecina (perturbacion gaussiana)
    new_x    = x + randn(1, 2);
    new_cost = ackley(new_x);

    % 2. Criterio de aceptacion
    if new_cost < best_cost          % Mejora -> aceptar siempre
        best_sol  = new_x;
        best_cost = new_cost;
        x = new_x;
    else
        delta_cost = new_cost - ackley(x);
        if rand < exp(-delta_cost / T) % Criterio de Metropolis
            x = new_x;
        end
    end

\end{lstlisting}
\end{frame}

\begin{frame}[fragile]{Lazo Principal del SA}
\begin{lstlisting}[language=Matlab, style=MATLABStyle]
    % 3. Registrar trayectoria y actualizar temperatura
    costs(iter) = best_cost;
    trace_x(end+1) = x(1);  trace_y(end+1) = x(2);
    trace_z(end+1) = ackley(x);
    T = T * alpha;           % Enfriamiento geometrico
    if T < T_min, break; end % Criterio de paro termico
end
\end{lstlisting}
\end{frame}
%------------------------------------------------
\section{Análisis de Resultados}
%------------------------------------------------

\begin{frame}{Visualización del Proceso de Búsqueda}
    \footnotesize
    \begin{columns}
        \column{0.5\textwidth}
        \begin{block}{Gráfica 3D: Superficie + Trayectoria}
            \begin{itemize}
                \item La superficie \texttt{surf()} muestra la función de Ackley
                \item Puntos rojos (\texttt{scatter3}) muestran la \textbf{trayectoria explorada}
                \item Se observa cómo al inicio el algoritmo salta ampliamente (alta $T$)
                \item Después converge hacia el mínimo
            \end{itemize}
        \end{block}
        \column{0.5\textwidth}
        \begin{block}{Gráfica 2D: Convergencia del costo}
            \begin{itemize}
                \item Eje $x$: número de iteración
                \item Eje $y$: $f_{\text{mejor}}$ en cada iteración
                \item Curva descendente confirma convergencia
                \item La curva no es monótonamente decreciente (SA acepta soluciones peores temporalmente)
            \end{itemize}
        \end{block}
    \end{columns}
    \vspace{0.2cm}
    % TODO: Add image Figures/SA_convergence.png
    \begin{exampleblock}{Salida esperada en consola}
        \texttt{Mejor solución encontrada: (0.0023, -0.0041)} \\
        \texttt{Costo de la mejor solución: 0.0123}
    \end{exampleblock}
\end{frame}

\begin{frame}{Resumen de Parámetros y su Efecto}
    \begin{table}[h]
        \centering
        \caption{Parámetros clave del Simulated Annealing}
        \scalebox{0.85}{
        \begin{tabular}{lccl}
            \toprule
            \textbf{Parámetro} & \textbf{Símbolo} & \textbf{Valor en ejemplo} & \textbf{Efecto} \\
            \midrule
            Temperatura inicial   & $T_0$     & 1000    & Mayor $\to$ más exploración inicial \\
            Temperatura mínima    & $T_{\min}$ & 1      & Criterio de paro del enfriamiento \\
            Factor de enfriamiento & $\alpha$  & 0.95   & Más cerca de 1 $\to$ enfriamiento lento \\
            Iteraciones máx.      & $N_{\max}$ & 700    & Límite de evaluaciones de la función \\
            Solución inicial      & $x_0$     & (0.1, 0.1) & Influye en calidad final \\
            \bottomrule
        \end{tabular}}
    \end{table}
\end{frame}

%------------------------------------------------
\section{Ventajas, Limitaciones y Variantes}
%------------------------------------------------

\begin{frame}{Ventajas y Limitaciones del SA}
    \footnotesize
    \begin{columns}
        \column{0.5\textwidth}
        \textbf{Ventajas:}
        \begin{itemize}
            \item \textbf{Generalidad}: aplica a cualquier problema con función de costo
            \item \textbf{Garantía asintótica}: con enfriamiento suficientemente lento, converge al óptimo global
            \item \textbf{Simplicidad de implementación}: pocos parámetros y código compacto
            \item \textbf{Robusto} ante ruido y discontinuidades
        \end{itemize}
        \column{0.5\textwidth}
        \textbf{Limitaciones:}
        \begin{itemize}
            \item \textbf{Sensibilidad al \textit{schedule}}: una mala elección de $\alpha$ o $T_0$ puede fallar
            \item \textbf{Convergencia lenta} para espacios de alta dimensión
            \item \textbf{No paralelizable} fácilmente (cadena de Markov secuencial)
            \item No garantiza el óptimo en tiempo finito
        \end{itemize}
    \end{columns}
    \vspace{0.3cm}
    \begin{alertblock}{Recomendación práctica}
        Para problemas de alta dimensión o multiobjetivo, considerar extensiones como \textbf{Parallel SA}, \textbf{MOSA} (Multi-Objective SA) o su combinación con metaheurísticas poblacionales.
    \end{alertblock}
\end{frame}

\begin{frame}{Variantes Avanzadas del SA}
    \begin{table}[h]
        \centering
        \caption{Comparación de variantes del Simulated Annealing}
        \scalebox{0.78}{
        \begin{tabular}{lll}
            \toprule
            \textbf{Variante} & \textbf{Característica principal} & \textbf{Aplicación} \\
            \midrule
            SA clásico         & Enfriamiento geométrico       & Optimización mono-objetivo \\
            Fast SA (FSA)      & Enfriamiento de Cauchy        & Convergencia más rápida \\
            Adaptive SA        & $T$ se ajusta dinámicamente   & Problemas dinámicos \\
            MOSA               & Frente de Pareto              & Optimización multiobjetivo \\
            Parallel SA        & Múltiples cadenas             & Cómputo paralelo / GPU \\
            SA + Híbrido       & SA + búsqueda local           & Optimización combinatoria \\
            \bottomrule
        \end{tabular}}
    \end{table}
\end{frame}

%------------------------------------------------
\section{Conclusiones}
%------------------------------------------------

\begin{frame}{Conclusiones}
    \begin{block}{Puntos clave}
        \begin{enumerate}
            \item El \textbf{Simulated Annealing} es una metaheurística inspirada en el proceso físico de recocido, capaz de escapar de mínimos locales mediante el criterio de \textbf{Metrópolis}.
            \item La temperatura controla el \textbf{balance exploración-explotación}: alta $T$ favorece diversidad; baja $T$ favorece convergencia.
            \item Los parámetros ($T_0$, $\alpha$, $T_{\min}$) deben elegirse con cuidado; su impacto es \textbf{determinante} para la calidad de la solución.
            \item Aplicado a la \textbf{función de Ackley}, el SA demuestra su capacidad de encontrar el mínimo global en un paisaje altamente multimodal.
        \end{enumerate}
    \end{block}
\end{frame}

\begin{frame}{Referencias}

    \begin{exampleblock}{Referencia fundamental}
        Kirkpatrick, S., Gelatt, C. D., \& Vecchi, M. P. (1983). Optimization by simulated annealing. \textit{Science}, 220(4598), 671--680.
        \url{https://doi.org/10.1126/science.220.4598.671}
    \end{exampleblock}
\end{frame}

\end{document}
